\documentclass{article}
\usepackage{amsmath,amsthm}
\usepackage{subcaption}
\usepackage{booktabs}
\usepackage{maleficium-interactive}
\maleficiumtheme{house}
\newtheorem{theorem}{Theorem}
\newtheorem{lemma}[theorem]{Lemma}
\theoremstyle{definition}
\newtheorem{definition}[theorem]{Definition}
\title{The Playground Paper}
\author{Maleficium}
\date{}
\begin{document}
\maketitle
\begin{abstract}
One source that compiles to a PDF today and is meant to become an
interactive HTML paper later. Every widget macro the package offers
appears at least once, beside live equations, plain figures, subfigures, math,
citations that preview their reference on hover, cross-references, and
theorems whose proofs fold away in the HTML paper.
\end{abstract}

\section{Introduction}
Text typeset with \TeX{}~\cite{knuth1984texbook} and \LaTeX{}~\cite{lamport1994latex}
(together,~\cite{knuth1984texbook,lamport1994latex}) reads the same in print and on screen; the widgets below add an
interactive layer on top of that~\cite{shannon1948mathematical}.
Section~\ref{sec:guide} explains each widget, and
Figures~\ref{fig:model}, \ref{fig:video}, \ref{fig:chart}, \ref{fig:plotly-tradeoff},
\ref{fig:plotly-latency} and~\ref{fig:html},
Table~\ref{tab:results} and Figures~\ref{fig:subs} and~\ref{fig:pair} show them.

\section{Widgets}\label{sec:widgets}

\subsection{Model}
\begin{figure}[ht]
\centering
\caption{A three-dimensional model widget; the PDF shows a render of the model.}
\label{fig:model}
\interactiveruntime[runtime=model@1, height=5cm, id=fig-model,
  alt={A colored cube, the interactive 3D mesh}]{models/mesh.glb}
\end{figure}

\subsection{Video}
\begin{figure}[ht]
\centering
\caption{A video widget; the PDF shows a frame of the clip.}
\label{fig:video}
\interactiveruntime[runtime=video@1, poster=figures/clip.png, height=5cm, id=fig-video,
  alt={A frame of the test-pattern clip}]{media/clip.mp4}
\end{figure}

\subsection{Table from CSV}
Table~\ref{tab:results} is read from a CSV file. The PDF keeps the first
three rows; the interactive view carries all five.
\begin{table}[ht]
\centering
\caption{Accuracy and latency for each configuration.}
\label{tab:results}
\interactiveruntime[runtime=table@1, pdfrows=3, id=tab-results,
  alt={Accuracy and latency for five configurations, from baseline to ours-xl}]{data/results.csv}
\end{table}

\subsection{Chart from a Vega-Lite spec}
\begin{figure}[ht]
\centering
\caption{Accuracy by method, drawn from a Vega-Lite specification.}
\label{fig:chart}
\interactiveruntime[runtime=chart@1, height=5cm, id=fig-chart,
  alt={Bar chart of accuracy rising from 0.71 for baseline to 0.89 for ours-xl}]{charts/ablation.vl.json}
\end{figure}

\subsection{Plotly figures from Python}
Figures~\ref{fig:plotly-tradeoff} and~\ref{fig:plotly-latency} are what
Python's \texttt{fig.write\_json()} saves, written from the same CSV by
\texttt{tools/plotly\_figures.py}. The first keeps Plotly Express defaults,
so it takes the paper's style. The second brings its own template, and
\texttt{template=figure} keeps it.
\begin{figure}[ht]
\centering
\caption{Accuracy against latency, a scatter plot in the paper's style.}
\label{fig:plotly-tradeoff}
\interactiveruntime[runtime=plotly@1, height=5cm, id=fig-plotly-tradeoff,
  alt={Line of accuracy against latency, from baseline at 12 ms to ours-xl at 58 ms}]{charts/tradeoff.plotly.json}
\end{figure}
\begin{figure}[ht]
\centering
\caption{Latency by method, a bar chart in the figure's own style.}
\label{fig:plotly-latency}
\interactiveruntime[runtime=plotly@1, template=figure, height=5cm, id=fig-plotly-latency,
  alt={Horizontal bars of latency, from 12 ms for baseline to 58 ms for ours-xl}]{charts/latency.plotly.json}
\end{figure}

\subsection{HTML widget}
\begin{figure}[ht]
\centering
\caption{A self-contained HTML widget in a sandboxed frame.}
\label{fig:html}
\interactive[height=5cm, id=fig-html,
  alt={The live HTML demo}]{widgets/demo/}
\end{figure}

\subsection{Widgets side by side}
Figure~\ref{fig:pair} holds one figure with two subfigures, each embedding
its own widget: the left panel reuses the chart spec as
Figure~\ref{fig:pair-a}, the right panel as Figure~\ref{fig:pair-b}.
\begin{figure}[ht]
\centering
\begin{subfigure}{0.45\textwidth}
\centering
\interactiveruntime[runtime=chart@1, height=3cm, id=fig-pair-a,
  alt={Small copy of the accuracy chart, left panel}]{charts/ablation.vl.json}
\caption{Left panel.}\label{fig:pair-a}
\end{subfigure}\hfill
\begin{subfigure}{0.45\textwidth}
\centering
\interactiveruntime[runtime=chart@1, height=3cm, id=fig-pair-b,
  alt={Small copy of the accuracy chart, right panel}]{charts/ablation.vl.json}
\caption{Right panel.}\label{fig:pair-b}
\end{subfigure}
\caption{Two widgets side by side.}
\label{fig:pair}
\end{figure}

\subsection{Widget in a float without a caption}
A float needs no caption. This one holds a widget and nothing else, so it
carries no label and no figure number.
\begin{figure}[ht]
\centering
\interactiveruntime[runtime=chart@1, height=3cm, id=chart-bare,
  alt={Small uncaptioned copy of the accuracy chart}]{charts/ablation.vl.json}
\end{figure}

A widget may also sit directly in the text, outside any float. It records
no label and no figure number either.

\begin{center}
\interactiveruntime[runtime=chart@1, height=3cm, id=chart-inline,
  alt={Small inline copy of the accuracy chart outside any float}]{charts/ablation.vl.json}
\end{center}

\section{Plain floats}
Not every figure is interactive. Figure~\ref{fig:plain} has a caption;
the one after it has none.
\begin{figure}[ht]
\centering
\includegraphics[height=3cm]{figures/plain.png}
\caption{A static figure with a caption.}
\label{fig:plain}
\end{figure}
\begin{figure}[ht]
\centering
\includegraphics[height=3cm]{figures/plain.png}
\end{figure}

\subsection{Subfloats}
Figure~\ref{fig:subs} uses the standard \texttt{subcaption} package, with
Figures~\ref{fig:suba} and~\ref{fig:subb} as its panels.
\begin{figure}[ht]
\centering
\begin{subfigure}{0.45\textwidth}
\centering
\includegraphics[width=\linewidth]{figures/left.png}
\caption{First panel.}\label{fig:suba}
\end{subfigure}\hfill
\begin{subfigure}{0.45\textwidth}
\centering
\includegraphics[width=\linewidth]{figures/right.png}
\caption{Second panel.}\label{fig:subb}
\end{subfigure}
\caption{Two subfigures side by side.}
\label{fig:subs}
\end{figure}

\section{Math and theorems}
Entropy is defined in Equation~\eqref{eq:entropy} and bounded in
Theorem~\ref{thm:bound}. In the HTML paper, hovering or focusing a
citation shows its reference, and each proof can be folded away.
\begin{definition}[Entropy~\cite{shannon1948mathematical}]\label{def:entropy}
The entropy of a random variable $X$ taking $n$ values with probabilities
$p_1, \dots, p_n$ is
\begin{equation}\label{eq:entropy}
H(X) = -\sum_{i=1}^{n} p_i \log p_i .
\end{equation}
\end{definition}
\begin{lemma}[Gibbs' inequality]\label{lem:gibbs}
For distributions $p$ and $q$ on the same $n$ outcomes,
$-\sum_i p_i \log p_i \le -\sum_i p_i \log q_i$~\cite[Thm.~2.6.3]{cover2006elements}.
\end{lemma}
\begin{proof}
Since $\log x \le x - 1$ for every $x > 0$, with equality only at $x = 1$,
\begin{align*}
\sum_i p_i \log \frac{q_i}{p_i}
  &\le \sum_i p_i \left( \frac{q_i}{p_i} - 1 \right) \\
  &= \sum_i q_i - \sum_i p_i \le 0 . \qedhere
\end{align*}
\end{proof}
\begin{theorem}\label{thm:bound}
For a distribution on $n$ outcomes, $0 \le H(X) \le \log n$.
\end{theorem}
\begin{proof}[Proof of Theorem~\ref{thm:bound}]
Each term of Equation~\eqref{eq:entropy} is nonnegative, so $H(X) \ge 0$.
Lemma~\ref{lem:gibbs} with the uniform $q_i = 1/n$ gives
$H(X) \le -\sum_i p_i \log \tfrac{1}{n} = \log n$.
\end{proof}

\subsection{Live equations}
A \verb|liveequation| is an equation whose parameters the reader can
change. The PDF shows the values at their defaults; the HTML paper puts a
slider and a number box under the equation and recomputes each result.
Equation~\eqref{eq:binary-entropy} is Definition~\ref{def:entropy} for two
outcomes with probabilities $p$ and $1-p$, in bits.
\begin{liveequation}[id=eq-binary-entropy,
    alt={Binary entropy in bits at a probability p the reader can change}]
  H_2(\maleficiumparam{p}{0.25}[0.01,0.99,0.01])
    = \liveresult[digits=3]{-(p*ln(p) + (1-p)*ln(1-p))/ln(2)}\ \text{bits}
  \label{eq:binary-entropy}
\end{liveequation}
A result may come before the parameters it uses, a parameter may be
shown again, and \verb|liveequation*| is unnumbered. With
\verb|mode=widget| the equation runs in a frame of its own, on the
figure plate, instead of in the article's text:
\begin{liveequation*}[id=eq-parabola, mode=widget,
    alt={A parabola y = a x squared with live a and x}]
  y = \liveresult{a*x^2}
    = \maleficiumparam{a}{1}[0.5,1.5,0.1] \cdot \maleficiumparam{x}{1.5}[-2,2,0.1]^2,
  \qquad x = \maleficiumparam{x}
\end{liveequation*}

\section{What each widget does}\label{sec:guide}
\begin{description}
\item[Model] \verb|runtime=model@1| loads a glTF model (\texttt{.glb}) for orbiting
  and zooming. The PDF shows the poster (Figure~\ref{fig:model}).
\item[Video] \verb|runtime=video@1| plays a clip, with an optional remote URL
  and checksum for large files. The PDF shows the poster frame (Figure~\ref{fig:video}).
\item[Table] \verb|runtime=table@1| reads a CSV into a sortable table. The PDF
  prints only the first \texttt{pdfrows} rows (Table~\ref{tab:results}).
\item[Chart] \verb|runtime=chart@1| renders a Vega-Lite specification with
  hover details; the PDF shows a poster (Figure~\ref{fig:chart}).
\item[Plotly] \verb|runtime=plotly@1| draws a figure saved from Python with
  Plotly's hover, in the paper's style unless \verb|template=figure|
  (Figures~\ref{fig:plotly-tradeoff} and~\ref{fig:plotly-latency}).
\item[HTML] \verb|\interactive| embeds a folder with an \texttt{index.html} in a
  sandboxed frame (Figure~\ref{fig:html}).
\item[Live equation] \verb|liveequation| with \verb|\maleficiumparam{name}{default}[min,max,step]|
  and \verb|\liveresult[digits=n]{formula}|: the PDF typesets the values at the
  defaults, the HTML paper lets the reader change them
  (Equation~\eqref{eq:binary-entropy}).
\end{description}
Every widget requires \texttt{alt=} text, which becomes the image
description in the HTML view. The app renders the PDF's picture of a
model, chart, Plotly figure or HTML widget from the widget itself, so these carry no
\texttt{poster=}; a video needs a \texttt{poster=} frame until video
posters are rendered.

\bibliographystyle{plain}
\bibliography{refs}
\end{document}
